\chapter{Data Description}
\label{ch:data_description}

This chapter outlines the two publicly available datasets that are used in this thesis, including their structure, class distributions and a comparison between the two. These two are from Kaggle and are examples of symptom-based disease prediction problems of varying sizes. The pipeline of pre-processing that was used for these datasets is described in Chapter~\ref{ch:methodology}.

\section{Dataset 1: Disease Symptom Prediction Dataset}
\label{sec:dataset1}

The first dataset (Disease Symptom Prediction Dataset~\citep{Kaggle_Disease2020}) includes 4,920 patient records, which label a set of symptoms for each of 41 different disease categories. Once the lists of symptoms are converted into binary vectors, there are 131 different symptoms in the set, each with a value of 1 if it is present and 0 if it is not. The raw file is well balanced with 120 records per disease, but the majority of these records are exact duplicates of each other, and once deduplicated only 304 unique symptom--disease combinations are left. The deduplication step (which is described in Chapter~\ref{ch:methodology}) is crucial because, if it were not done, the training and testing on duplicated rows would lead to an optimistic estimation of performance.

The main features of this data are summarized in Table~\ref{tab:dataset1_overview}.

\begin{table}[H]
\centering
\caption{Overview of Dataset~1: Disease Symptom Prediction Dataset}
\label{tab:dataset1_overview}
\begin{tabular}{@{}ll@{}}
\toprule
\textbf{Property}         & \textbf{Value} \\
\midrule
Source                     & Kaggle (itachi9604) \\
Number of records          & 4,920 raw (304 unique after deduplication) \\
Number of features         & 131 (binary symptoms) \\
Number of disease classes  & 41 \\
Feature type               & Binary (0/1) \\
Target variable            & Disease name (categorical) \\
Missing values             & Present in raw data \\
\bottomrule
\end{tabular}
\end{table}

The data set was collected from a number of medical resources along with some extra data like disease descriptions, symptom intensities and recommended precautions. I only employ the symptom-disease mapping because I want to test the performance of the classifiers on the main prediction task.

\subsubsection{Class Distribution}

When the data from Dataset~1 is presented in its raw form, this set is perfectly balanced; there are 120 records for each of the 41 diseases. But this balance is artificial, since most of the 120 records are identical. Following deduplication, each disease has between 5 and 10 different symptom--disease combinations; this results in a very slight imbalance (2:1) between the largest and smallest classes.

This is a residual imbalance which is small compared to that of Dataset~2, but is still significant enough to be dealt with to make sure that the classifiers do not slightly favour the better-represented diseases. For this reason, class weighting is used throughout (Section~\ref{sec:class_imbalance}).

\subsubsection{Symptom Overlap Among Diseases}

An interesting feature of this set of data is that a number of the diseases have similar symptoms. There are many different categories of illness that present with non-specific symptoms, such as fever: this is medically reasonable, because there are a lot of illnesses that can have fever as a symptom, for example flu and tuberculosis. The difficulty for the classifier is then to identify the set of symptoms, instead of just one symptom.

A review of the data reveals that certain combinations of symptoms are very good predictors of a disease group. For instance, the mixture of yellowish skin, loss of appetite and abdominal pain suggests liver ailments, and skin rash with itching and a nodular skin eruption suggests dermatological ailments. These joint patterns can be represented by tree-based classifiers like Random Forest or XGBoost, which can split on multiple features at once.

\section{Dataset~2: Diseases and Symptoms Dataset}
\label{sec:dataset2}

The second is the Diseases and Symptoms Dataset~\citep{Kaggle_Diseases2023}, which is considerably larger, with approximately 246,000 patient records and significantly more conditions than Dataset~1. Following removal of disease classes with insufficient samples for stratifying the train--test split, and deduplication of identical records (Chapter~\ref{ch:methodology}), approximately 189,600 records are available across 727 disease classes for the experiments.

Table~\ref{tab:dataset2_overview} provides an overview of its main features.

\begin{table}[H]
\centering
\caption{Overview of Dataset~2: Diseases and Symptoms Dataset}
\label{tab:dataset2_overview}
\begin{tabular}{@{}ll@{}}
\toprule
\textbf{Property}         & \textbf{Value} \\
\midrule
Source                     & Kaggle (dhivyeshrk) \\
Number of records          & $\sim$246,000 raw ($\sim$189,600 used) \\
Number of features         & 377 binary symptoms \\
Number of disease classes  & 727 (after filtering rare classes) \\
Feature type               & Binary (0/1) \\
Target variable            & Disease name (categorical) \\
Missing values             & Minimal \\
\bottomrule
\end{tabular}
\end{table}

The larger size of Dataset~2 allows me to observe how classifiers scale with additional data. With almost 50 times more records than Dataset~1, this dataset produces more robust performance estimates and a lower risk of overfitting. It also decreases the variability in the estimates derived from cross-validation, thus enhancing the reliability of the comparison.

\subsubsection{Class Distribution}

In contrast to Dataset~1, Dataset~2 is truly imbalanced: the most common diseases show up with many records, whereas the rarest diseases only show up a few times. In some cases, classes were too small to have a stratified train--test split and were omitted during preprocessing, leaving 727 classes. This is important because some classifiers (in particular Naive Bayes) require a sufficient number of examples per class for the probability estimates to be reasonable.

The common ones are generally more general and are often very common diseases, while the rare ones are more specific diseases. The residual imbalance even after filtering is significant, but not quite as severe as in other clinical datasets found in the literature, in which the maximum-to-minimum number of class instances can be thousands to one.

\subsubsection{Data Quality}

Large web-sourced datasets are known to have an issue in terms of data quality. I sampled 100 of the records in Dataset~2, and made sure that the symptom-disease associations were sensible. Most were, but there were some records that seemed to be missing some symptoms. For instance, one of the records included just two symptoms for a diabetes diagnosis, even though diabetes typically presents with more symptoms.

Although there are some quality problems here and there, the size of the dataset allows a small amount of noisy records to have little impact on the performance of the classifiers. The benefit of having 246{,}000 examples outweighs the cost of slightly noisier labels compared to a smaller, more carefully curated dataset.

\section{Dataset Comparison}
\label{sec:dataset_comparison}

In addition to size, there are a number of important differences between the two datasets. Table~\ref{tab:dataset_comparison} summarizes the key differences.

\begin{table}[H]
\centering
\caption{Comparison of Dataset~1 and Dataset~2}
\label{tab:dataset_comparison}
\begin{tabular}{@{}lll@{}}
\toprule
\textbf{Characteristic} & \textbf{Dataset~1} & \textbf{Dataset~2} \\
\midrule
Number of records & 4,920 (304 unique) & $\sim$246,000 ($\sim$189,600 used) \\
Number of classes & 41 & 727 \\
Number of features & 131 & 377 \\
Class imbalance ratio & $\sim$2:1 (after dedup) & High \\
Missing values & Moderate & Minimal \\
Data quality & High (curated) & Good (web-sourced) \\
Smallest class size & 5 examples & $\geq$2 examples \\
\bottomrule
\end{tabular}
\end{table}

Dataset~1 is small and clean, but with a concern of overfitting for larger models with many parameters. For instance, XGBoost with hundreds of trees can memorize training examples instead of learning the general pattern. Dataset~2 is noisier than Dataset~1, as it was collected in a different way, but has more information. Larger training sets typically result in more powerful classifiers and less variance in performance estimates.

I selected two sets of data of very different scale, to not only be able to compare the classifiers in absolute terms, but also how they scale with data availability. There are classifiers that fare poorly on a small data set but do better as the data set grows larger. Naive Bayes is frequently used as a baseline on small datasets, because it requires few examples to estimate class probabilities. Its performance tends to flatten out as the dataset becomes larger, as the independence assumption restricts the amount of learning that the model can do. A second advantage of ensemble methods like XGBoost is that they can continuously improve as they receive more information, since they can model more complicated decision boundaries.

This is important for deployment considerations, as all healthcare settings have different training data sets. Patients can be routinely handled by a large hospital with tens of thousands of historical records in its care, while a small clinic may not have anywhere near that many. If a classifier only works on a large set of data, then it is not appropriate to use in low-resource settings.

With two data sets it is also possible to test the stability of the ranking of classifiers. If Random Forest outperforms SVM on Dataset~1 but SVM outperforms Random Forest on Dataset~2, the ranking is dataset-dependent rather than indicative of any general superiority. Otherwise, if Random Forest gives better accuracy than SVM on both datasets, that is more convincing evidence that it is really more appropriate to this task.
